\input pdfToolbox
\input pdfmsym
\pdfmsymsetscalefactor{10}
\input ../preamble
\def\enumindent{0cm}

\def\printmcount{\the\counter{section}.\the\counter{math counter}}
\setcounter{section}{4}

\centerline{\setfontandscale{bf}{20pt}Commutative Algebra \& Algebraic Geometry}
\medskip
\centerline{\bf Exercise \the\counter{section}}
\smallskip
\centerline{\it Ari Feiglin}

\bprob

    The {\emphcolor prime avoidance lemma} says that if $\p_1,\dots,\p_n\in\spec(R)$ are prime ideals and $I\leq R$ is
    an ideal contained in $\bigcup_i\p_i$ then $I\leq\p_i$ for some $i$.
    Induct on $n$ to prove the contrapositive: if $I$ is not contained in any $\p_i$ then it is not contained in the union.
    \benum
        \item Prove the base case.
        \item For $n>1$ use induction to find for any $i$ an $x_i\in I-\p_1\cup\cdots\cup\hat\p_i\cup\cdots\cup\p_n$.
        \item Show that either $x_i\notin\p_i$ for some $i$, or
        $$ \sum_i x_1\cdots\hat x_i\cdots x_n \in I-\bigcup_i\p_i , $$
        and conclude the lemma.
        \item Prove the similar result: if $\p\in\spec(R)$ is prime and $I_1,\dots,I_n\leq R$ are ideals such that
        $\bigcap_iI_i\leq\p$ then $I_i\leq\p$ for some $i$.
    \eenum

\eprob

\benum
    \item The base case is trivial, since the union is just $\p_1$.

    \item Since $I$ is not contained in $\p_i$ for all $i$, applying the inductive hypothesis to the $n-1$ prime ideals
    $\p_1,\dots,\hat\p_i,\dots,\p_n$ shows that $I$ is not contained in $\bigcup_{j\neq i}\p_j$.
    Thus there is an $x_i\in I$ not in $\bigcup_{j\neq i}\p_j$, as desired.

    \item Let $y_i=x_1\cdots\hat x_i\cdots x_n$.
    Assume the claim is not true: $x_i\in\p_i$ for all $i$ and $\sum_iy_i\in\p_j$ for some $j$.
    Now for $i\neq j$, $y_i$ is in $\p_j$ since it is a product of $x_j\in\p_j$.
    Since the sum of all $y_i$s is in $\p_j$, so too must $y_j$ be, so $y_j\in\p_j$.
    $\p_j$ is a prime ideal, so $x_i\in\p_j$ for some $i\neq j$.
    But then $x_i\in\bigcup_{j\neq i}\p_j$, a contradiction.

    In either case we have an element in $I-\bigcup_i\p_i$, so $I$ is not contained in the union of the $\p_i$s.

    \item Suppose not, so $I_i$ is not contained in $\p$ for all $i$.
    Let $x_i\in I_i-\p$, and define $y=x_1\cdots x_n$.
    But then, since $I_i$ is an ideal, $y\in I_i$ for all $i$ so $y\in\bigcap_iI_i$ and thus $y\in\p$.
    Since $\p$ is prime, $x_i\in\p$ for some $i$, a contradiction.
\eenum

\bprob

    \benum
        \item Find an example of a ring which is Noetherian but not Artinian.
        \item Let $0\to K\to M\to L\to0$ be a SES of $R$-modules.
        Show that $M$ is Artinian iff $K$ and $L$ are.
        \item Let $M$ be an Artinian $R$-module and $S\subseteq R$ a multiplicatively closed set.
        Show that $S^{-1}M$ is Artinian.
        \item Let $R=K$ be a field and $V$ a vector space over $K$.
        Show that $V$ is finite-dimensional iff it is both Artinian and Noetherian.
        \item Assume $R$ is Artinian.
        \benum
            \item Let $p\in\spec(R)$.
            Show that $R/\p$ is an Artinian domain.
            Furthermore let $x\in R/\p$ be nonzero; show that there is a $y\in R/\p$ and $n\in\bN$ such that
            $x^n=x^{n+1}y$.
            Conclude that $xy=1$ and $\p$ is maximal.
            \item Show that $R$ has a finite number of maximal ideals.
            \item Conclude that $\spec(R)$ is a finite discrete space.
        \eenum
    \eenum

\eprob

\benum
    \item The ring of integers $\bZ$ is Noetherian but not Artinian.
    It is Noetherian because $(n)\subseteq(m)$ iff $m$ divides $n$.
    Thus an ideal is only contained in a finite number of ideals: the ideals generated by its divisors.
    Since integers have finitely many divisors, this is finitely many.
    In particular every ascending chain of ideals must stabilize.

    But it is not Artinian since for $n\neq\pm1$, the sequence $I_k=(n^k)$ is an infinite descending chain of ideals.
    It is descending because $n^k$ divides $n^{k+1}$ so $I_{k+1}\subseteq I_k$, and $n^k\in I_k-I_{k+1}$.

    \item Let $0\to K\xto fM\xto gL\to0$ be a SES.
    Suppose that $M$ is Artinian; we will show that $K$ and $L$ are.
    If $K_n\subseteq K,L_n\subseteq L$ are descending chains of submodules, then $f(K_n),g^{-1}(L_n)\subseteq M$
    are as well.
    By $M$'s Artinian-ness, they must stabilize, so $f(K_n)=f(K_{n+1})$ and $g^{-1}(L_n)=g^{-1}(L_{n+1})$ eventually.
    By $f$'s injectivity and $g$'s surjectivity, this means $K_n=K_{n+1}$ and $L_n=L_{n+1}$ eventually
    (the direct image of an injective map is injective, and the inverse image of a surjective map is injective).
    Thus $K,L$ are Artinian, as all their descending chains of submodules stabilize.

    Now suppose that $K,L$ are Artinian, we will show that $M$ is.
    Let $M_n\subseteq M$ be a descending chain of submodules.
    Then $f^{-1}(M_n)\subseteq K,g(M_n)\subseteq L$ are descending chains of submodules of $K,L$ respectively,
    and thus must stabilize.
    So eventually $f^{-1}(M_n)=f^{-1}(M_{n+1})$ and $g(M_n)=g(M_{n+1})$ eventually, say for all $n\geq N$.
    We claim that this means $M_n=M_{n+1}$; the proof of this is the same as in lecture.

    Since $M_{n+1}\subseteq M_n$ we need only show the inverse inclusion.
    Let $x\in M_n$, so $g(x)\in g(M_n)=g(M_{n+1})$ so that $g(x)=g(y)$ for some $y\in M_{n+1}$.
    Then $x-y\in\ker g=\im f$ by short exactness, so $x-y=f(z)$ for $z\in K$.
    Since $x-y\in M_n$ this means that $z\in f^{-1}(M_n)=f^{-1}(M_{n+1})$, we see that $f(z)\in M_{n+1}$.
    Thus $f(z)=x-y\in M_{n+1}$ and since $y$ is already in $M_{n+1}$ this means $x\in M_{n+1}$ as desired.

    Thus all descending chains of submodules of $M$ stabilize, so it is Artinian.

    \item If $V$ is finite-dimensional, then all of its chains of submodules are bound by its dimension.
    Indeed, if $0=V_0\subsetneq V_1\subsetneq\cdots\subsetneq V_n=V$ is a chain of submodules, considering dimensions
    we must have $\dim V_i<\dim V_{i+1}$.
    Thus we must have $n\leq\dim V$, and since this is finite, $V$ has finite length and is thus Artinian and Noetherian.

    Conversely if $V$ is Artinian and Noetherian, it has finite length $n$.
    Take a chain of length $n$ (a composition series) $0=V_0\subsetneq\cdots\subsetneq V_n=V$.
    The only simple $K$-modules are those of dimension $1$, so $1=\dim V_{i+1}/V_i=\dim V_{i+1}-\dim V_i$.
    Thus
    $$ \dim V = \dim V_n = \sum_i \dim V_{i+1} - \dim V_i = n $$
    so $V$ is finite-dimensional.

    \item
    \benum
        \item Point (2) shows that quotients of Artinian modules are Artinian, so $R/\p$ is Artinian.
        Since $\p$ is prime, as usual $R/\p$ is a domain.

        Consider the chain of ideals $I_n=(x^n)$.
        These are descending since $x^{n+1}=xx^n$ and so $x^{n+1}\in I_n$ implying $I_{n+1}\subseteq I_n$.
        Since $R/\p$ is Artinian, this means that $I_n=I_{n+1}$ eventually.
        Thus $x^n\in(x^{n+1})$ and so there is a $y\in R/\p$ such that $x^n=yx^{n+1}$ as desired.
        This means that $x^n(xy-1)=0$ and since $R/\p$ is an integral domain either $x^n=0$ or $xy=1$.
        Again since $R/\p$ is a domain and $x\neq0$, $x^n\neq0$, so $xy=1$.

        So every nonzero element of $R/\p$ has an inverse, meaning $R/\p$ is a field.
        Thus $\p$ is maximal.

        \item Let $\Sigma$ be the set consisting of finite intersections of maximal ideals of $R$:
        $$ \Sigma = \set{\bigcap_{i=1}^n\m_i}[\m_i\in\specm(R)] . $$
        Since $R$ is Artinian, its poset of ideals has DCC and thus $\Sigma$ has a minimal element.
        Let this be $I=\bigcap_{i=1}^n\m_i$.
        For any maximal ideal $\m\in\specm(R)$, since $\m\in\Sigma$, by minimality $I\leq\m$.
        By our first question, since $\m_i$ are prime, this means that $\m_i\leq\m$ for some $i$.
        By maximality, $\m_i=\m$, meaning that $\m_1,\dots,\m_n$ exhaust all of the maximal ideals of $R$.

        \item Since all of $R$'s prime ideals are maximal, $R$ has finitely many prime ideals, so $\spec(R)$ is finite.
        Furthermore as we showed in a previous homework, maximal ideals are closed points.
        So all singletons in $\spec(R)$ are closed, and thus all subsets of $\spec(R)$ are closed (all subsets are finite
        and thus a finite union of closed sets, so closed).
        Thus $\spec(R)$ is discrete.
    \eenum
\eenum

\bprob

    Let $\phi\colon R\to S$ be a ring morphism.
    \benum
        \item Let $I\leq R$ be an ideal.
        \benum
            \item Show that there is a unique morphism $\bar\phi\colon R/I\to S/IS$ making the following diagram commute:

            \commsq{R&S\cr R/I&S/IS}
                {\phi}{\bar\phi}{}{}

            (the vertical arrows are the canonical projections.)

            \item We showed previously that $\spec(R/I)$ and $\spec(S/IS)$ can be identified with $V_R(I)$ and
            $V_S(IS)$ respectively.
            Show that under this identification, $\bar\phi^*\colon V_S(IS)\to V_R(I)$ is the restriction of
            $\phi^*\colon\spec(S)\to\spec(R)$ to $V_S(IS)$.
        \eenum

        \item Let $\p\in\spec(R)$.
        Show that $(\phi^*)^{-1}(\p)$ is homeomorphic to $\spec(S_\p/\p S_\p)$.
    \eenum

\eprob

\benum
    \item
    \benum
        \item Let $\pi_R,\pi_S$ be the associated projections.
        We need
        $$ \bar\phi\circ\pi_R = \pi_S\circ\phi . $$
        If $\bar\phi'$ is another morphism satisfying this, then $\bar\phi\circ\pi_R=\bar\phi'\circ\pi_R$.
        Since $\pi_R$ is surjective, and thus epic, this implies $\bar\phi=\bar\phi'$, showing uniqueness.

        Now $\phi(I)\subseteq IS$ by definition, so $I\subseteq\ker(\pi_S\circ\phi)$, meaning we can quotient
        $\pi_S\circ\phi$ by $I$.
        That is, $\bar\phi$ exists.
        Explicitly,
        $$ \bar\phi(x+I) = \phi(x) + IS , $$
        and the above argument shows that this is well-defined.

        \item The identification of $\spec(R/I)$ with $V_R(I)$ is given by the topological embedding induced by the canonical
        projection, $\pi_R^*\colon\spec(R/I)\to\spec(R)$.
        The commutative diagram

        \commsq{R&S\cr R/I&S/IS}
            {\phi}{\bar\phi}{\pi_R}{\pi_S}

        Induces a commutative diagram of affine schemes (by functorality):

        \commsq{\spec(S/IS)&\spec(R/I)\cr S&R}
            {\bar\phi^*}{\phi^*}{\pi_S^*}{\pi_R^*}

        Meaning that
        $$ \pi_R^*\circ\bar\phi^* = \phi^*\circ\pi_S^* . $$
        Since $\pi_R$ is a topological embedding, we can take its inverse and
        $$ \bar\phi^* = (\pi_R^*)^{-1}\circ\phi^*\circ\pi_S^* , $$
        as desired.
    \eenum

    \item We know that
    $$ (\phi^*)^{-1}(\p) = \set{\q\in\spec(S)}[\phi^*(\q)=\p] = \set{\q\in\spec(S)}[\p=\phi^{-1}(\q)] . $$
    Furthermore we know that $\spec(S_\p/\p S_\p)$ is homeomorphic to $V_{S_\p}(\p S_\p)$, which is
    $$ V_{S_\p}(\p S_\p) = \set{\q S_\p}[\q\in\spec(S),\q\cap\phi(S-\p)=\emptyset,\p S_\p\leq\q S_\p] , $$
    since prime ideals of $S_\p$ are of the form $\q S_\p$ for $\q\in\spec(S)$ not intersecting the localizing set.

    Now consider the continuous map induced by the structure morphism $\iota^*\colon\spec(S_\p)\to\spec(S)$.
    We know from homework 2 that $\iota^*$ is a topological embedding, so that $V_{S_\p}$ is homeomorphic
    to its image under $\iota^*$.

    Since $\iota^*(\q S_\p)=\q$ for prime ideals $\q\in\spec(S)$, its image is
    $$ \iota^* V_{S_\p} = \set{\q\in\spec(S)}[\q\cap\phi(S-\p)=\emptyset,\p S_\p\leq\q S_\p] . $$
    Now we claim that this is equal to $(\phi^*)^{-1}(\p)$: so $\p S_\p\leq\q S_\p$ and $\q\cap\phi(S-\p)=\emptyset$ iff
    $\p=\phi^{-1}(\q)$.
    If $\p=\phi^{-1}(\q)$ then $\phi(\p)\leq\q$ and thus $\p S_\p\leq\q S_\p$, and if $x\notin\p$ then $\phi(x)\notin\q$
    so that $\q\cap\phi(S-\p)=\emptyset$.

    Conversely if $\p S_\p\leq\q S_\p$ then $\phi(\p)S_\p\leq\q S_\p$ and since $\iota^{-1}$ preserves order, we have
    $$ \phi(\p)S_\p \leq \q S_\p \implies \iota^{-1}(\phi(\p)S_\p) \leq \iota^{-1}(\q S_\p) = \q . $$
    Since $\phi(\p)\leq\iota^{-1}(\phi(\p)S_\p)$ we have $\phi(\p)\leq\q$ and thus $\p\leq\phi^{-1}\q$.
    And $\q\cap\phi(S-\p)=\emptyset$ implies that if $\phi(x)\in\q$ then $\phi(x)\notin\phi(S-\p)$ so $x\in\p$.
    Thus $\phi^{-1}\q\leq\p$, meaning $\phi^{-1}\q=\p$.

    So we have that
    $$ \spec(S_\p/\p S_\p) \cong V_{S_\p}(\p S_\p) \cong \iota^* V_{S_\p}(\p S_\p) = (\phi^*)^{-1}(\p) , $$
    as desired.
\eenum

\bprob

    Let $\phi\colon R\to A$ be a ring morphism.
    \benum
        \item Let $R=\bC[t]$ and $A=\bC[t,t^{1/2},t^{1/3},\dots]$ and $\phi$ be inclusion.
        Show that $\phi$ is integral but not finite.
        \item Suppose $\phi$ is integral, and let $J\leq A$ be an ideal.
        Show that the induced morphism $\bar\phi\colon R/\phi^{-1}J\to A/J$ is injective and integral.
        \item Suppose $\phi$ is integral, show that $\phi^*\colon\spec(A)\to\spec(R)$ is closed.
        \item Suppose $\phi$ is finite, and let $\p\in\spec(R)$.
        Show that $(\phi^*)^{-1}(\p)$ is homeomorphic to a finite discrete set.
    \eenum

\eprob

\benum
    \item Since the integral closure is a subalgebra, it is sufficient to show that $t^{k/n}$ is integral over $\bC[t]$
    for all $k,n$.
    This is clear, though since $x^n-t^k\in R[x]$ is a monic polynomial with root $t^{k/n}$.
    Thus $A$ is integral over $R$.

    But $A$ is not finite: for any finite subset $B\subset A$, let $k_1/n_1,\dots,k_\ell/n_\ell$ be all the exponents
    of $t$ occuring in $B$, for $k_i,n_i$ coprime.
    Then
    $$ R[B] \subseteq R[\set{t^{k_i/n_i}}[i=1,\dots,\ell]] \qquad \hbox{as $R$-modules.} $$
    For $t^k\in R$ notice that $t^k t^{k_i/n_i}=t^{(k_i+kn_i)/n_i}$ does not increase the denominator.
    Thus for any $p(t)\in R$, the denominators of the exponents in $p(t)t^{k_i/n_i}$ are bound by $n_i$.
    Thus all the exponents in $B$ are of the form $k/n$ where $n$ is bound by $\max n_i$.
    Taking $N$ to be greater than all $n_i$, we see that $t^{1/N}$ cannot be in $R[B]$, so $R[B]\neq A$.
    So $A$ is not finitely-generated as an $R$-module, and thus not a finite $R$-algebra.

    \item This was shown in lecture, nevertheless we repeat it here.
    Note that the kernel of $\phi$ composed with the projection $A\to A/J$ has kernel $\phi^{-1}J$, so the quotient
    $\bar\phi$ is injective (the kernel of $f\circ g$ is $(f\circ g)^{-1}1=g^{-1}\ker f$).
    Then for $\bar a\in R/\phi^{-1}J$ and $\bar x\in A/J$, $\bar a\bar x=\bar\phi(\bar a)\bar x$ which is just
    $\overline{\phi(a)x}$.

    For $\bar x\in A/J$, lift it to $x\in A$ which is integral over $R$, so there is a monic polynomial such that
    $x^n+a_1x^{n-1}+\cdots+a_n=0$.
    Then let $\bar a_i$ be the projections of $a_i$, so
    $$ \bar x^n + \bar a_1\bar x^{n-1} + \cdots \bar a_n = 0 $$
    and thus $\bar x$ is integral over $A/\phi^{-1}J$.

    \item We need to show that for any ideal $J\leq A$, $\phi^*V_A(J)$ is closed in $\spec(R)$, so is a vanishing
    set itself.
    By the above point $\bar\phi\colon R/\phi^{-1}J\to A/J$ is injective and integral.
    In lecture we showed that this means $\bar\phi^*$ is surjective.
    Now consider the commutative diagram

    \commsq{\spec(A/J)&\spec(R/\phi^{-1}J)\cr\spec(A)&\spec(R)}
        {\bar\phi^*}{\phi^*}{\pi_A^*}{\pi_R^*}

    so
    $$ \pi_R^*\circ\bar\phi^* = \phi^*\circ\pi_A^* . $$
    Recall that $\pi_A^*$ is a topological embedding whose image is $V_A(J)$.
    Thus
    $$ \phi^*V_A(J) = \pi_R^*\bar\phi^*\spec(A/J) = \pi_R^*\spec(R/\phi^{-1}J) = V_R(\phi^{-1}J) , $$
    so $\phi^*V_A(J)$ is indeed closed, as desired.

    \item From the previous question we know $(\phi^*)^{-1}(\p)$ is homeomorphic to $\spec(A_\p/\p A_\p)$.
    Since $A$ is a finite $R$-algebra, $A_\p$ is a finite $R_\p$-algebra.
    This is because if $R^n\to A$ is a surjective $R$-module morphism, then by exactness of localization,
    $R^n_\p\to A_\p$ is a surjective $R_\p$-module morphism, so $A_\p$ is finitely generated as an $R^n_\p$-module,
    thus a finite $R_\p$-algebra.

    If we denote by $\phi_\p\colon R_\p\to A_\p$ the localization of $\phi$, then $\phi_\p^{-1}(\p A_\p)=\p R_\p$.
    $K=R_\p/\p R_\p$ is a field, and $A_\p/\p A_\p$ is a $K$-algebra.
    We claim that it is a finite $K$-algebra.

    In general if $A$ is a finite $R$-algebra, and $J\leq A$ is an ideal with preimage $J_R\leq R$,
    then $A/J$ is a finite $R/J_R$-algebra.
    This is because if $a_1,\dots,a_n\in A$ generate it as an $R$-module, then their projections
    $\bar a_1,\dots,\bar a_n\in A/J$ generate it as an $R/J_R$-module, since multiplication by $R/J_R$ is just projection
    of multiplication by $R$.

    Being a finite $K$-module means having finite dimension, and as we showed in question 2 this implies being Artinian.
    Thus $A_\p/\p A_\p$ is an Artinian $K$-module, and thus an Artinian ring (its ideals are submodules).
    Again from question 2, we showed that this means $\spec(A_\p/\p A_\p)$ is a finite discrete space.
    Since $(\phi^*)^{-1}(\p)$ is homeomorphic to it, it is also a finite discrete space.
\eenum

\bprob

    \benum
        \item Prove Noetherian induction.
        \item Let $X$ be a topological space.
        Prove that the following are equivalent:
        \benum
            \item $X$ is Noetherian,
            \item Every subspace of $X$ is quasi-compact,
            \item Every open subspace of $X$ is quasi-compact.
        \eenum
        \item Let $R=\bC[x_1,x_2,\dots]/I$ where $I=(\set{x_ix_j}[i,j\geq1])$.
        \benum
            \item Prove that $R$ is not Noetherian.
            \item Prove that $\spec(R)$ consists of a single point, and is thus Noetherian.
        \eenum
    \eenum

\eprob

\benum
    \item Suppose $\Sigma$ is a poset with DCC, and $P$ is a boolean map out of $\Sigma$.
    Further suppose that $P(p)=1$ when $p$ is minimal, and for non-minimal $p$ such that $P(q)=1$ for all $q<p$, then
    $P(p)=1$.
    We want to show that $P$ is identically $1$ on $\Sigma$.

    Suppose not, so the set
    $$ \Sigma' = \set{p\in P}[P(p)=0] $$
    is nonempty.
    Because $\Sigma$ has DCC, $\Sigma'$ has a minimal element, let it be $p\in\Sigma'$.
    Because $P(p)=0$, by assumption $p$ cannot be minimal in $\Sigma$.
    Furthermore for every $q<p$, by $p$'s minimality $q\notin\Sigma'$ so $P(q)=1$.
    Then by assumption $P(p)=1$, a contradiction.

    \item Note that $X$ is Noetherian iff its poset of open subsets has ACC.
    This is because the poset of open subsets of $X$ is reverse order-isomorphic to $X$'s poset of closed subsets
    via taking complements.

    If $X$ is Noetherian and $S\subseteq X$ is any subspace, we will show that $S$ is quasi-compact.
    Let $\set{U_i}$ be an open cover of $S$, and let $\Sigma$ be the set of all finite unions of $U_i$s.
    $\Sigma$ consists of open subsets of $X$ and thus has a maximal element, say $\bigcup_{i=1}^nU_i$ (ignoring issues
    with indexing, the point is clear).
    This must be an open cover of $S$, as if $x\in S$ is not in the union, then since $\set{U_i}$ covers $S$ there must
    be some $j$ such that $x\in U_j$.
    But then $\bigcup_{i=1}^nU_i\cup U_j$ is a proper superset of $\bigcup_{i=1}^nU_i$, and a finite union of $U_i$s,
    contradicting $\bigcup_{i=1}^nU_i$'s maximality.

    Clearly if every subspace of $X$ is quasi-compact, then every open subspace of $X$ is quasi-compact.

    Now suppose that every open subspace of $X$ is quasi-compact; we will show $X$ is Noetherian by showing its poset
    of open subspaces has ACC.
    Let $U_1\subseteq U_2\subseteq\cdots$ be an increasing chain of open subspaces of $X$.
    Then $\set{U_i}_i$ forms an open cover of its union $U=\bigcup_iU_i$.
    Since by assumption $U$ is quasi-compact, there is a finite subcover $U_{i_1},\dots,U_{i_n}$ so $U=\bigcup_{j=1}^nU_{i_j}$.
    Without loss of generality $i_1<\cdots<i_n$, and so $U_{i_1}\subseteq\cdots\subseteq U_{i_n}$ since the $U_i$
    form an increasing chain.
    Thus
    $$ U = \bigcup_{j=1}^nU_{i_j} = U_{i_n} . $$
    Then for every $i\geq i_n$, $U_i\subseteq U$ by definition, and so $U_i\subseteq U_{i_n}$ and so $U_i=U_{i_n}$
    (since the chain is increasing).
    So the chain stabilizes at $i_n$.
    Thus the poset of open subspaces of $X$ has ACC, so $X$ is Noetherian.

    \item Ideals of $R$ are in order-preserving bijection with ideals of $\bC[x_1,x_2,\dots]$ containing $I$.
    So we show that there is an infinite ascending chain of ideals of $\bC[x_1,x_2,\dots]$ containing $I$.
    Define
    $$ I_n = (\set{x_1,\dots,x_n}\cup\set{x_ix_j}[i,j>n]) , $$
    so $I_0=I$ and $I_n\subseteq I_{n+1}$.
    But $x_{n+1}\in I_{n+1}-I_n$, so this inclusion is strict.
    Thus $R$ is not Noetherian.

    The prime ideals of $R$ are in bijection with the prime ideals of $\bC[x_1,x_2,\dots]$ containing $I$.
    Suppose $I\leq\p$ for a prime ideal $\p$.
    Then for all $i$, $x_i^2\in I\subseteq\p$ so by primality $x_i\in\p$.
    So $(x_1,x_2,\dots)\leq\p$, but also $\m=(x_1,x_2,\dots)$ is a maximal ideal.
    This is because the map $\bC[x_1,x_2,\dots]\to\bC$ given by mapping a polynomial to its constant coefficient is
    a ring morphism and surjective.
    Its kernel is $\m$: clearly $\m$ is contained in the kernel, and if
    $$ f = \sum_\alpha a_\alpha x^\alpha $$
    has constant coefficient zero, then each $\alpha$ (of which there are finitely many) has some nonzero index.
    For each $\alpha$ let $i_\alpha$ be a nonzero index and $\alpha_i$ be obtained by removing this index from $\alpha$.
    Then
    $$ f = \sum_\alpha a_\alpha x_{i_\alpha}x^{\alpha_i} \in \m . $$
    Since $\m$ is maximal, this means that $\m=\p$.
    So $\m$ is the only prime ideal containing $I$, thus $\spec(R)=\set\m$, as desired.
\eenum

\bprob

    Let $R$ be a ring.
    \benum
        \item Show that if $R$ is Noetherian, it has finitely many minimal prime ideals.
        \item Assume $R$ is Noetherian and let $I\leq R$ be an ideal.
        Show that $I$ contains a finite product of prime ideals.
        \item Let $R\bbracks x$ be the $R$-algebra of power series with coefficients in $R$.
        Let $f=\sum_{n=0}^\infty a_nx^n\in R\bbracks x$.
        Show that if $R$ is Noetherian, then $f$ is nilpotent iff each $a_n$ is nilpotent.
        \item Give a counterexample to the above point when $R$ is not Noetherian.
        \item Let $M$ be a Noetherian $R$-module.
        Show that $R/\ann(M)$ is a Noetherian ring.
        \item Let $I_1,\dots,I_n\leq R$ be ideals such that $\bigcap_{i=1}^nI_i=0$ and $R/I_i$ is Noetherian for each $i$.
        Show that $R$ is Noetherian.
    \eenum

\eprob

\benum
    \item If $R$ is Noetherian, then $\spec(R)$ is a Noetherian space, and thus as we showed in recitation has
    finitely many irreducible components.
    In the previous homework we showed that an irreducible component of $\spec(R)$ is the same as the vanishing set
    of a minimal prime ideal.
    Since $V$ is injective on prime ideals (directly by definition), this means that $R$ has finitely many minimal prime
    ideals, as desired.

    \item Assume there is an ideal which does not contain a finite product of prime ideals.
    Since $R$ is Noetherian, there is a maximal ideal which does not contain a finite product of prime ideals, let this
    be $I$.
    $I$ itself cannot be prime, as a prime ideal contains itself and thus a finite product of prime ideals.
    So there are $x,y\in R$ such that $x,y\notin I$ but $xy\in I$.
    Let $J_1=(I,x)$ and $J_2=(I,y)$; these are proper extensions of $I$ and thus each contains a finite product
    of prime ideals, $P_1\subseteq J_1$ and $P_2\subseteq J_2$.
    Notice though that $J_1J_2=(I^2,xI,yI,xy)\subseteq I$, and thus $P_1P_2\subseteq J_1J_2\subseteq I$.
    $P_1P_2$ is a finite product of prime ideals, so $I$ does contain a finite product of prime ideals, a contradiction.

    \item Suppose that $f$ is nilpotent; we proceed the same as in homework 1.
    Let $\p$ be a prime ideal of $R$, and consider the projection $\pi\colon R\to R/\p$ which extends to
    $\bar\pi\colon R\bracks x\to(R/\p)\bracks x$ point-wise.
    Since nilpotent elements are preserved under morphisms, $\bar\pi(f)$ is nilpotent.
    But since $R/\p$ is a domain, so is $(R/\p)\bracks x$, and thus $\bar\pi(f)=0$.

    Indeed, if $A$ is an integral domain then $A\bracks x$ is as well: if $f=\sum_na_nx^n$ is nonzero and $g=\sum_nb_nx^n$
    is a power series such that $fg=0$, then we will show inductively that $b_n=0$.
    Let $a_m$ be the minimal nonzero coefficient in $f$, and let $c_n$ be the coefficients of $fg$, so $c_n=0$.
    Then
    $$ c_m = \sum_{i=0}^m a_ib_{m-i} = a_mb_0 $$
    and since $a_m\neq0$, we have $b_0=0$.
    Now inductively assume $b_i=0$ for all $i<k$, then
    $$ 0 = c_{m+k} = \sum_{i=0}^{m+k} b_ia_{m+k-i} = \sum_{i=k}^{m+k}b_ia_{m+k-i} = \sum_{i=0}^m a_ib_{m+k-i} = a_mb_k , $$
    and thus $b_k=0$.
    So $g=0$, and thus we have shown $A\bracks x$ is a domain.

    Returning to our question $\bar\pi(f)=0$ so $a_n\in\p$ for all $n$.
    This is true for all prime $\p$ and thus $a_n\in\nil(R)$ as $\nil(R)$ is the intersection of all prime ideals.
    So if $f$ is nilpotent, so are all the $a_n$s.
    Note that this is true for all rings, not just Noetherian ones.

    For the other direction we must somehow use Noetherian-ness.
    Suppose $f=\sum_na_nx^n$ where $a_n$ is nilpotent for each $n$.
    Let $I_k=(a_0,\dots,a_k)$ so that $I_k$ is an ascending chain of ideals.
    Since $R$ is Noetherian, eventually this stabilizes, say at $\ell$, so $I_k=I_\ell$ for all $k\geq\ell$.
    Let $N_i$ be the order of nilpotency of $a_i$, so $a_i^{N_i}$, and let $N=N_0+\cdots+N_\ell$.

    We claim that $f^N=0$.
    Let $c_n^k$ be the coefficient of $x_n$ in $f^k$.
    Then
    $$ c_n^{k+1} = \sum_{i=0}^n c_n^ka_{n-i} . $$
    Inductively we will show that $c_n^k\in I_\ell^k$.

    For $k=1$ then $c_n^1=a_n$, which is in $I_\ell$ (since $I_n\subseteq I_\ell$).
    Now by induction,
    $$ c_n^{k+1} = \sum_{i=0}^n c_n^ka_{n-i} \in \sum_{i=0}^n a_{n-i}I_\ell^k \subseteq \sum I_\ell I_\ell^k = I_\ell^{k+1} . $$
    As desired.

    Now we claim that $I_\ell^N=0$.
    $I_\ell^N$ is generated by elements of the form $a=\prod_{i=0}^\ell a_i^{k_i}$ where $\sum k_i=N$.
    Then for some $i$ we must have $k_i\geq N_i$, as otherwise $\sum k_i<\sum N_i=N$ a contradiction.
    But then $a_i^{k_i}=0$ so $a=0$.
    Thus $I_\ell^N=0$, and in particular $c_n^N=0$ for all $n$.
    Then $f^N=0$, so $f$ is nilpotent.

    \item We showed that $f$ being nilpotent means all of its coefficients are in any ring, so we need to give an example
    of a power series whose coefficients are nilpotent but $f$ is not.
    Let us consider
    $$ R = \bF_2[t,t^{1/2},t^{1/3},\dots]/(t) , $$
    and
    $$ f = \sum_{n=1}^\infty t^{1/n}x^n . $$
    Each $t^{1/n}$ is nilpotent as raising it to the $n$th power gives $t$, which is zero in $R$.
    
    Now notice that for $g=\sum_{n=0}^\infty a_nx^n$ where $a_n$ are taken from an $\bF_2$-algebra,
    $g^2=\sum a_n^2x^{2n}$.
    If we let $b_n$ be the coefficients of $g^2$ then
    $$ b_n = \sum_{k=0}^n a_ka_{n-k} . $$
    Since if $k\neq n-k$, $a_ka_{n-k}$ appears twice in the sum (once when $k=k$ and once when $k=n-k$), it cancels out
    and so we can focus only on values of $k$ where $k=n-k$.
    That is, $2k=n$, and so
    $$ b_{2k} = a_k^2 $$
    and all other values of $b_n$ are zero.
    Thus
    $$ g^2 = \sum b_nx^n = \sum b_{2k}x^{2k} = \sum a_k^2x^{2k} $$
    as desired.

    Inductively we see then that
    $$ g^{2^k} = \sum a_n^{2^k}x^{2^kn} . $$
    In our case
    $$ f^{2^k} = \sum_{n=1}^\infty t^{2^k/n}x^{2^kn} , $$
    which is nonzero since $t^{2^k/n}$ is nonzero for $n>2^k$.
    Thus $f$ is not nilpotent.

    \item We first claim that if $R$ has a faithful Noetherian module $M$, then $R$ is Noetherian.
    $M$ being faithful means that the defining morphism $\phi\colon R\to\endo(M)$ is an injection.
    Thus if we show $\endo(M)$ is a Noetherian $R$-module, then we will be finished as submodules of Noetherian
    modules are Noetherian.

    Since $M$ is Noetherian, it is finitely generated, say by $m_1,\dots,m_n\in M$.
    Then the surjection $f\colon R^n\to M$ mapping $e_i\mapsto m_i$ defines a map
    $$ \endo(M) = \hom(M,M) \to \hom(R^n,M) \cong \hom(R,M)^n \cong M^n $$
    given by $g\mapsto g\circ f$.
    Since $f$ is epic, the induced map is an injection.
    Finite direct products of Noetherian modules are Noetherian, so $M^n$ is Noetherian.
    Submodules of Noetherian modules are Noetherian, so $\endo(M)$ is Noetherian as desired.

    And $M$ is a faithful $R/\ann(M)$-module, and also a Noetherian $R/\ann(M)$-module (since the submodules of $M$
    as an $R$-module vs as an $R/\ann(M)$-module are the same).
    So by what we just showed, $R/\ann(M)$ is Noetherian.

    \item Let $\pi_i\colon R\to R/I_i$ be the canonical projection maps.
    Define the map
    $$ \Pi\colon R \to \prod_{i=1}^n R/I_i $$
    whose $i$th coordinate is $\pi_i$, i.e.\ $\Pi(r)=(\pi_i(r))_i$.
    This is an $R$-module morphism since each $\pi_i$ is (and the universal property of direct products).
    Furthermore $\Pi$ is injective, since $\Pi(r)=0$ iff $\pi_i(r)=0$ for all $i$ iff $r\in\bigcap_iI_i=0$.

    The product of all $R/I_i$ is Noetherian, since each $R/I_i$ is Noetherian and the property of being Noetherian is
    preserved under finite products.
    Thus $R$ can be embedded into a Noetherian module, and is therefore Noetherian itself.
\eenum

\def\T{{\c T}}

\bprob

    Let $V$ be a finite-dimensional complex vector space.
    Let $\T\subseteq\endo_\bC(V)$ be a commuting set of diagonalizable endomorphisms.
    Show that there exists a basis of $V$ such that each vector in it is an eigenvector of every $T\in\T$.

\eprob

Let $T$ be a diagonalizable linear operator and $U$ a $T$-invariant subspace.
Let $V_{\lambda_1},\dots,V_{\lambda_n}$ be $T$'s eigenspaces so that
$$ V = V_{\lambda_1} \oplus \cdots \oplus V_{\lambda_n} . $$
We claim that
$$ U = \bigoplus_{i=1}^n(V_{\lambda_i}\cap U) . $$
Right-to-left inclusion is clear.
Now for $u\in U$, $u$ is equal to $v_1+\cdots+v_n$ for $v_i\in V_{\lambda_i}$; we claim that $v_i\in U$ as well.
Note that $T$ quotients to a linear operator $\bar T\colon V/U\to V/U$ since $U$ is $T$-invariant, and $\bar v_i$, the projections of $v_i$, are eigenvectors
of $\bar T$ with respective eigenvalues $\lambda_i$.
Since $\lambda_i$ are distinct so $\bar v_i$ are linearly independent, or zero.
But $v_1+\cdots+v_n\in U$ so $\bar v_1+\cdots+\bar v_n=0$, and thus $\bar v_i$ must all be zero, i.e.\ $v_i\in U$ as desired.

Note that if $T,S$ commute and $V_\lambda$ is an eigenspace of $T$, then $V_\lambda$ is $S$-invariant.
Indeed, for $v\in V_\lambda$,
$$ TSv = STv = S(\lambda v) = \lambda Sv , $$
so $Sv\in V_\lambda$.

Now say we have finitely many linear operators $T_1,\dots,T_n$ which commute and are diagonalizable.
Let $\lambda^i_j$ be the eigenvalues of $T_i$.
Then $V_{\lambda^i_j}$ is $T_k$-invariant for all $i,j,k$.
Thus inductively we see
$$ V = \bigoplus_{j_1,\dots,j_n} \bigcap_i V_{\lambda^i_{j_i}} . $$
For $n=1$ this is clear, and then inductively $\bigcap_iV_{\lambda^i_{j_i}}$ is $T_{n+1}$-invariant so as we showed
$$ \bigcap_{i=1}^nV_{\lambda^i_{j_i}} = \bigoplus_{j_{n+1}}\parens{\bigcap_{i=1}^n V_{\lambda^i_{j_i}}\cap V_{\lambda^{n+1}_{j_{n+1}}}} . $$
Thus
$$ V = \bigoplus_j\bigcap_i V_{\lambda^i_j} = \bigoplus_{j_1,\dots,j_{n+1}}\bigcap_i V_{\lambda^i_{j_i}} $$
as desired.

For each $J=(j_1,\dots,j_n)$, let $V_J=\bigcap_iV_{\lambda^i_{j_i}}$, so $V=\bigoplus_JV_J$.
Notice that vectors in $V_J$ are eigenvectors of each $T_i$ (of eigenvalue $\lambda^i_{j_i}$).
Thus if we take a basis $B_J$ for each $V_J$, their union $B$ forms a basis for $V$, and elements of this basis are eigenvectors for all $T_i$.
Thus $B$ diagonalizes each $T_i$.

Now in our case consider the subspace spanned by $\T$, which still consists of commuting endomorphisms.
Since $\T$ is a subset of $\endo(V)$, this is finite-dimensional, and thus can be finitely generated by a finite subset of $\T$.
So let $T_1,\dots,T_n\in\T$ span this subspace.
As we showed, they are simultaneously diagonalizable, say by a basis $B$.
But then every endomorphism in their span is diagonalized by $B$ (since the linear combination of diagonal matrices is diagonal).
In particular every endomorphism in $\T$ is diagonalized by $B$, so $\T$ is simultaneously diagonalizable.

\bye

